\section{Core Equations}

This section introduces the minimal equation system used throughout the paper. The equations are not tied to a specific biological system. They provide a general form for describing state change, flow, and stable structure.

\subsection{Basic Variables}

The basic fields are the information density
\begin{equation}
  \I(\mathbf{r},t)
\end{equation}
and the information flow
\begin{equation}
  \J(\mathbf{r},t).
\end{equation}
The field \(\I\) represents local density of structure or constraint. The field \(\J\) represents the direction and rate of state change.

\subsection{Continuity Equation}

Information density obeys the continuity equation
\begin{equation}
  \frac{\partial \I}{\partial t} + \nabla \cdot \J = \sigma,
  \label{eq:continuity}
\end{equation}
where \(\sigma(\mathbf{r},t)\) is a local generation rate. This equation records local production, outflow, and conservation-like transport of the coarse-grained quantity \(\I\).

\subsection{Structure of the Flow}

The flow is written as
\begin{equation}
  \J = -D\nabla \I + \mu \mathbf{F},
  \label{eq:flow}
\end{equation}
where \(D\) is a diffusion coefficient, \(\mu\) is a response coefficient, and \(\mathbf{F}\) is an induced direction field.

\subsection{Potential}

The direction field is determined by the induced potential:
\begin{equation}
  \mathbf{F} = -\nabla \PhiI.
  \label{eq:force}
\end{equation}
The potential is a functional of information density:
\begin{equation}
  \PhiI = \PhiI[\I].
  \label{eq:potential}
\end{equation}
The specific form of \(\PhiI\) is system-dependent and left unspecified.
Thus, biases in information density induce a gradient-driven bias in state change.
The sign convention is fixed only after \(\PhiI[\I]\) is specified.
For example, choosing \(\PhiI\sim-\I\) makes the induced component
\(-\nabla\PhiI\) point toward increasing constraint density.
Other biological observables may require the opposite convention or additional
response terms.

\subsection{Dynamical Loop}

Equations~\eqref{eq:continuity}--\eqref{eq:potential} form a loop: the distribution of \(\I\) induces \(\PhiI\), the potential shapes \(\J\), and the flow redistributes \(\I\). In general this loop does not necessarily terminate in a static equilibrium. It can produce non-equilibrium structures that are sustained under continued update.

\subsection{Two-Layer Description of Time}

The equations use the external time parameter \(t\). An internal time \(\tau\) can be introduced as
\begin{equation}
  \frac{\dd \tau}{\dd t} = \rho(\I,\J),
\end{equation}
where \(\rho(\I,\J)\) is a state-dependent progression rate. The variable \(t\) is used for equations of motion. The variable \(\tau\) is used to interpret the internal progress of the system.

\subsection{Stable Structures}

Under these dynamics, a region of state space is stable if states within or near it tend to remain there or return after perturbation. Such a region is interpreted as an attractor.

\subsection{Summary}

The minimal equation system is
\begin{align}
  \frac{\partial \I}{\partial t} + \nabla \cdot \J &= \sigma,\\
  \J &= -D\nabla \I + \mu \mathbf{F},\\
  \mathbf{F} &= -\nabla \PhiI,\\
  \PhiI &= \PhiI[\I].
\end{align}
This system describes state change, flow formation, and the emergence of stable structures in a common form.
